% !TeX root = ../papers/nc-p3.tex
\section{正則性座標の水平切片と添字保持同型曲線}
\label{sec:rank-fifteen-horizontal-fibers}
\begin{theorem}\label{thm:rank-fifteen-horizontal-fibers}
標数零の代数閉体で $y=c\ne0$ を固定する。階数15のPauli商曲面の添字保持同型類の
水平切片は、$t\ne0$ 上の曲線
\[
       V^2=2t\bigl((-t^2+2t-4c)^3-216ct^3\bigr)
\]
である。$c\ne1,1/64$ では滑らかな完備超楕円曲線の二つの
Weierstrass点を除いたもので、種数は3、全点の最小周期は2である。
一位置移動の商は種数2の曲線 $W^2=(x^3-27c)((1-x)^2-4c)$ から
無限遠点を除いたものである。
ある有限集合 $E\subset\overline{\mathbb Q}^{\times}$ の外では
水平切片の全点が標準Koszul左右AS正則で、中心四次式上の自由階数は各始点で256である。
\end{theorem}
\begin{proof}
定理\ref{thm:rank-fifteen-rational-quotient}、第\ref{sec:rank-fifteen-global-invariants}節および
定理\ref{thm:rank-fifteen-formal-boundary}より、$p^2=x^3-27c$, $q^2=(1-x)^2-4c$
の各点が一つの添字保持同型類を与え、一位置移動は $(p,q)\mapsto(-p,-q)$ である。
$t=1-x+q$ と取れば $(1-x+q)(1-x-q)=4c$ より
\[
 x=1-(t+4c/t)/2,\qquad q=(t-4c/t)/2,\qquad V=4t^2p .
\]
従って表示はアフィン曲線の同型である。一位置移動は
$(t,V)\mapsto(4c/t,-16c^2V/t^4)$ となる。
$0,\infty$ はいずれも二重被覆の分岐点であり、この二点だけを除く。
移動不変環は $k[x,pq]$ である。実際、$k[x]$ 上の自由基底
$1,p,q,pq$ のうち $1,pq$ が固定されるので $W=pq$ と取れる。

二つの被開平式の終結式は $(c-1)^2(1-64c)$ である。
\path{computations/verify_rank_fifteen_horizontal_fibers.py} は $t$ 表示の六次因子の
判別式が非零定数と $c^{10}(c-1)^2(64c-1)$ の積であることを検査する。
六次因子の定数項も非零なので、通常の切片の二重被覆は八点で分岐し種数3、
移動の商は五次式の五根と無限遠で分岐し種数2となる。
同時に $p=q=0$ となる点が無いので、最小周期は2である。
定理\ref{thm:rank-fifteen-vertical-curves}の有限例外集合の $y$ 座標を $E$ とすれば
全次数の主張が従い、通常表示外の四点は定理\ref{thm:rank-fifteen-formal-boundary}で覆われる。
\end{proof}
\begin{corollary}
$c=1,1/64$ の水平切片にはそれぞれ一つの通常二重点があり、
切片の正規化の種数は2、移動の商の正規化の種数は1である。
各二重点が表す窓は最小周期1の標準Koszul左右AS正則窓であり、
その切片の他の全点は最小周期2である。
\end{corollary}
\begin{proof}
二重点の $(x,t)$ は順に $(3,-2),(3/4,1/4)$ である。
七次式にある唯一の二重根を $b$ とすると、$(t-b)^2$ を取り除いた五次式は平方自由で、
$b$ で非零である。よって二つの異なる接線があり、正規化の二点が一つの同型類に重なる。
商の五次式も唯一の二重因子を除くと平方自由三次式となる。
二重点は $p=q=0$ の二つの既知の正則窓である。
\end{proof}
特に $c=3$ は滑らかな種数3切片である。3の上で $\nu(3)=1$ とすれば、
この曲線の $\nu(x)>1/3$ を満たす全点は定理\ref{thm:rank-fifteen-overlap-annulus}により
全次数で認証される。この明示的な範囲には、例えば $x\in3\mathbb Z_3$ の全ての持上げが入る。
